\documentclass[a4paper,12pt]{article}
\usepackage[T2A]{fontenc}
\usepackage[utf8]{inputenc}
\usepackage[russian]{babel}
\usepackage{amsmath,amssymb,mathtools}
\usepackage{helvet}
\renewcommand{\familydefault}{\sfdefault}
\usepackage{icomma}
\usepackage[margin=18mm]{geometry}
\usepackage{microtype}
\usepackage{tikz}
\usetikzlibrary{calc,arrows.meta,positioning,shapes.geometric}
\usepackage{tabularx,booktabs,longtable,array}
\usepackage{enumitem}
\usepackage{parskip}
\usepackage{fancyhdr}
\usepackage{setspace}
\usepackage{xcolor}
\usepackage[hidelinks]{hyperref}

\definecolor{accent}{HTML}{1F6F78}
\definecolor{darkaccent}{HTML}{174F56}
\definecolor{textgray}{HTML}{2F3437}
\definecolor{softgray}{HTML}{F3F5F5}

\setlength{\parindent}{0pt}
\onehalfspacing
\pagestyle{fancy}
\setlength{\headheight}{14pt}
\fancyhf{}
\fancyhead[L]{\small\color{textgray}Текстовые задачи: смеси и проценты}
\fancyhead[R]{\small\color{textgray}29.09.26}
\fancyfoot[C]{\small\thepage}
\setlist[itemize]{leftmargin=*,itemsep=0.25em,topsep=0.35em}
\setlist[enumerate]{leftmargin=*,itemsep=0.45em,topsep=0.35em}

\newcommand{\sectiontitle}[1]{\vspace{0.4em}{\Large\bfseries\color{accent}#1}\par\vspace{0.35em}}
\newcommand{\key}[1]{\textbf{\color{darkaccent}#1}}
\newcommand{\formula}[1]{\[\boxed{#1}\]}

\begin{document}

\begin{titlepage}
\thispagestyle{empty}
\centering
\vspace*{2.3cm}
{\Large\color{textgray}Чек-лист\par}
\vspace{0.7cm}
{\Huge\bfseries\color{accent}Текстовые задачи:\par}
\vspace{0.15cm}
{\Huge\bfseries\color{accent}смеси, растворы и проценты\par}
\vspace{1.3cm}
{\large Екатерина Гнедкова\par}
\vspace{0.4cm}
{\large 29.09.26\par}
\vfill
{\large\color{textgray}Лёвин Артём Александрович эксклюзивно для Екатерины Гнедковой\par}
\vspace{0.9cm}
\begin{tikzpicture}[remember picture,overlay]
  \draw[accent!65,line width=1.2pt]
    ([xshift=15mm,yshift=14mm]current page.south west)
    .. controls +(45mm,18mm) and +(-40mm,-5mm) ..
    ([xshift=-15mm,yshift=20mm]current page.south east);
\end{tikzpicture}
\end{titlepage}

\sectiontitle{1. Главная идея: одна таблица для разных сюжетов}
Раствор, сплав, смесь и даже задача про высушивание продукта удобно рассматривать как систему из \key{целого} и выбранного \key{компонента}. В большинстве задач достаточно трёх величин:

\begin{center}
\renewcommand{\arraystretch}{1.35}
\begin{tabularx}{0.92\linewidth}{>{\bfseries}l X}
\toprule
Величина & Смысл \\
\midrule
Масса смеси & масса всего раствора, сплава, продукта и т.п. \\
Концентрация & доля выбранного компонента в целом \\
Масса компонента & сколько выбранного вещества фактически содержится в смеси \\
\bottomrule
\end{tabularx}
\end{center}

Ключевая связь:
\formula{m_{\text{компонента}}=m_{\text{смеси}}\cdot c}
где $c$ записана \key{десятичной дробью}: $15\%=0{,}15$, $41\%=0{,}41$.

\sectiontitle{2. Что можно складывать при смешивании}
Если смешивают несколько частей без потерь вещества, то складываются массы:
\[
 m_{\text{итог}}=m_1+m_2+\dots
\]
и массы выбранного компонента:
\[
 m_{\text{комп, итог}}=m_{\text{комп},1}+m_{\text{комп},2}+\dots
\]

\key{Концентрации напрямую не складываются.} Итоговая концентрация определяется через отношение массы компонента к общей массе смеси.

\sectiontitle{3. Базовый пример: смешивание растворов}
Смешали $4$ кг $15\%$-го раствора и $6$ кг $25\%$-го раствора.

\begin{center}
\renewcommand{\arraystretch}{1.35}
\begin{tabular}{lccc}
\toprule
 & Масса смеси & Концентрация & Масса вещества \\
\midrule
1-й раствор & $4$ & $0{,}15$ & $0{,}6$ \\
2-й раствор & $6$ & $0{,}25$ & $1{,}5$ \\
Итог & $10$ & $x$ & $10x$ \\
\bottomrule
\end{tabular}
\end{center}

Уравнение:
\[
10x=0{,}6+1{,}5=2{,}1,
\qquad x=0{,}21=21\%.
\]

\key{Самопроверка:} итоговая концентрация должна лежать между $15\%$ и $25\%$.

\sectiontitle{4. Чистая вода и чистое вещество}
\begin{itemize}
\item Для чистой воды (если отслеживаем растворённое вещество): $c=0$.
\item Для чистого вещества: $c=1$.
\item Добавление воды меняет общую массу, но не меняет массу растворённого вещества.
\item Добавление чистого вещества меняет и общую массу, и массу компонента на одну и ту же добавленную величину.
\end{itemize}

\sectiontitle{5. Сплавы решаются тем же методом}
Для сплава роль «растворённого вещества» играет выбранный металл. Например, если ищем никель, то
\[
m_{\text{никеля}}=m_{\text{сплава}}\cdot c_{\text{никеля}}.
\]
Если смешали $x$ кг $10\%$-го и $y$ кг $30\%$-го сплава и получили $200$ кг $25\%$-го сплава, то:
\[
\begin{cases}
x+y=200,\\
0{,}1x+0{,}3y=50.
\end{cases}
\]

\sectiontitle{6. Высушивание: ищите неизменную часть}
При высушивании уходит вода, а масса сухого вещества сохраняется. Поэтому полезно выбрать именно \key{сухое вещество} как отслеживаемый компонент.

Пример: изюм массой $20$ кг содержит $95\%$ сухого вещества, исходный виноград --- $10\%$ сухого вещества. Если $x$ --- масса винограда, то
\[
0{,}10x=20\cdot0{,}95=19,
\qquad x=190\text{ кг}.
\]

\key{Главный вопрос:} какая величина остаётся неизменной между начальным и конечным состоянием?

\clearpage
\thispagestyle{fancy}
\begin{center}
{\LARGE\bfseries\color{accent}Алгоритм: задача на смесь или раствор}
\end{center}
\vspace{0.9cm}

\begin{center}
\begin{tikzpicture}[
  node distance=10mm and 14mm,
  startstop/.style={ellipse,draw=accent,very thick,align=center,minimum width=43mm,minimum height=10mm,fill=softgray},
  process/.style={rounded corners=2mm,draw=accent,very thick,align=center,text width=64mm,minimum height=11mm},
  decision/.style={diamond,aspect=2.0,draw=accent,very thick,align=center,text width=40mm,inner sep=2pt},
  arrow/.style={-{Stealth[length=3mm]},very thick,draw=darkaccent}
]
\node[startstop] (start) {Прочитайте условие};
\node[process,below=of start] (component) {Выберите компонент, массу которого удобно отслеживать};
\node[process,below=of component] (table) {Составьте таблицу: масса смеси / концентрация / масса компонента};
\node[decision,below=12mm of table] (cases) {Сколько неизвестных?};
\node[process,text width=55mm,below left=15mm and 7mm of cases] (one) {Одна неизвестная: составьте одно уравнение по массе компонента};
\node[process,text width=55mm,below right=15mm and 7mm of cases] (two) {Две неизвестные: добавьте уравнение по общей массе или второму состоянию};
\node[process,below=22mm of cases] (solve) {Решите уравнение или систему};
\node[process,below=of solve] (check) {Проверьте смысл ответа: диапазон концентрации, знак, единицы};
\node[startstop,below=of check] (finish) {Запишите ответ};

\draw[arrow] (start)--(component);
\draw[arrow] (component)--(table);
\draw[arrow] (table)--(cases);
\draw[arrow] (cases)--node[above left]{1}(one);
\draw[arrow] (cases)--node[above right]{2}(two);
\draw[arrow] (one.south) |- (solve.west);
\draw[arrow] (two.south) |- (solve.east);
\draw[arrow] (solve)--(check);
\draw[arrow] (check)--(finish);
\end{tikzpicture}
\end{center}
\clearpage

\sectiontitle{7. Процентные задачи: исходное значение и коэффициент изменения}
Процент удобно сразу переводить в коэффициент.

\begin{center}
\renewcommand{\arraystretch}{1.35}
\begin{tabular}{lll}
\toprule
Изменение & Коэффициент & Формула \\
\midrule
увеличили на $p\%$ & $1+\frac{p}{100}$ & $B=A\left(1+\frac{p}{100}\right)$ \\
уменьшили на $p\%$ & $1-\frac{p}{100}$ & $B=A\left(1-\frac{p}{100}\right)$ \\
\bottomrule
\end{tabular}
\end{center}

Пример: после снижения цены на $13\%$ товар стоит $1218$ руб. Тогда
\[
0{,}87x=1218,
\qquad x=1400\text{ руб.}
\]

\key{Типичная ошибка:} после снижения на $13\%$ использовать коэффициент $1{,}13$. Для снижения нужен коэффициент $0{,}87$.

\sectiontitle{8. Несколько последовательных изменений}
Последовательные проценты \key{не складываются}; перемножаются коэффициенты.

Например, увеличение на $15\%$, затем снижение на $20\%$:
\[
1{,}15\cdot0{,}80=0{,}92.
\]
Итоговая величина составляет $92\%$ от исходной, то есть уменьшилась на $8\%$.

\sectiontitle{9. Быстрые ориентиры перед решением}
\begin{itemize}
\item Если смешивают --- складывайте массы смеси и массы выбранного компонента.
\item Если добавляют воду --- масса компонента сохраняется.
\item Если выпаривают воду --- масса сухого/растворённого вещества сохраняется.
\item Если задача про сплав --- выберите один металл и отслеживайте его массу.
\item Если есть два разных состояния одной смеси --- часто нужны две таблицы и система уравнений.
\item Если цена изменилась на процент --- сначала определите, повышение это или снижение, затем используйте коэффициент.
\end{itemize}

\sectiontitle{10. Частые ошибки}
\begin{itemize}
\item складывать концентрации вместо масс вещества;
\item писать $15$ вместо $0{,}15$ в формуле массы компонента;
\item забывать, что концентрация воды равна $0$, а чистого вещества --- $1$;
\item выбирать компонент, масса которого меняется, хотя в задаче есть неизменная часть;
\item путать увеличение и снижение цены;
\item складывать последовательные проценты вместо умножения коэффициентов;
\item не проверять, попадает ли концентрация смеси между исходными концентрациями.
\end{itemize}

\clearpage
\sectiontitle{Тренировочный блок — Путь А}
\textbf{10 задач.} Решайте через таблицу там, где это уместно. В тренировочном блоке решения не приводятся.

\begin{enumerate}
\item Смешали $8$ кг $10\%$-го раствора и $12$ кг $30\%$-го раствора. Найдите концентрацию полученной смеси.

\item К $12$ кг $30\%$-го раствора добавили $3$ кг воды. Найдите новую концентрацию.

\item К $18$ кг $20\%$-го раствора добавили $2$ кг чистого вещества. Найдите концентрацию полученного раствора.

\item Смешали $x$ кг $10\%$-го сплава и $y$ кг $30\%$-го сплава. Получили $200$ кг сплава с содержанием выбранного металла $25\%$. Найдите массы исходных сплавов.

\item Из винограда, содержащего $10\%$ сухого вещества, получили $20$ кг изюма, содержащего $95\%$ сухого вещества. Сколько килограммов винограда использовали?

\item В первом опыте смешали $x$ кг $40\%$-го раствора, $y$ кг $70\%$-го раствора и $20$ кг воды и получили раствор концентрации $41\%$. Во втором опыте вместо воды взяли $20$ кг $60\%$-го раствора, а итоговая концентрация стала $53\%$. Найдите $x$ и $y$.

\item В первом сосуде концентрация кислоты равна $x$, во втором --- $y$. Смешали $5$ кг первого и $10$ кг второго раствора и получили $15$ кг $40\%$-го раствора. Затем смешали $10$ кг первого и $5$ кг второго раствора и получили $15$ кг $50\%$-го раствора. Найдите концентрации исходных растворов и массу кислоты в первых $5$ кг первого раствора.

\item После снижения цены на $13\%$ товар стал стоить $1218$ руб. Найдите первоначальную цену.

\item Цену товара сначала увеличили на $15\%$, затем уменьшили на $20\%$. После двух изменений товар стоит $1840$ руб. Найдите первоначальную цену.

\item Имеется $30$ кг $20\%$-го раствора. Сколько килограммов воды нужно выпарить, чтобы концентрация стала $30\%$?
\end{enumerate}

\clearpage
\sectiontitle{Ответы}
\begin{center}
\renewcommand{\arraystretch}{1.5}
\begin{tabularx}{0.9\linewidth}{>{\bfseries}c X}
\toprule
№ & Ответ \\
\midrule
1 & $22\%$ \\
2 & $24\%$ \\
3 & $28\%$ \\
4 & $50$ кг и $150$ кг \\
5 & $190$ кг \\
6 & $x=50$ кг, $y=30$ кг \\
7 & $60\%$ и $30\%$; масса кислоты в $5$ кг первого раствора --- $3$ кг \\
8 & $1400$ руб. \\
9 & $2000$ руб. \\
10 & $10$ кг \\
\bottomrule
\end{tabularx}
\end{center}

\vfill
\begin{center}
\small\color{textgray}Лёвин Артём Александрович эксклюзивно для Екатерины Гнедковой
\end{center}

\end{document}